\chapter{Migration}
\label{ch:migration}

This chapter provides a complete migration guide for teams moving from v1 DRI runbooks to
gert v2 + \texttt{gert.ops}. It covers the automated migration tool, manual migration
patterns, and a validation checklist.

% ─────────────────────────────────────────────────────────────────────────────
\section{Migration Overview}
\label{sec:migration-overview}
% ─────────────────────────────────────────────────────────────────────────────

Migrating to \texttt{gert.ops} involves two independent concerns:

\begin{description}
  \item[Core v1 → v2 migration] Schema field promotions, step type renames, and new required
    fields. Handled by \texttt{gert migrate}. See the \emph{gert v2 Design Document}, §11 for
    the core migration reference.
  \item[\texttt{gert.ops} Kit adoption] Adding the Kit declaration, replacing informal manual
    steps with typed \texttt{ops.*} steps, declaring roles, adding approval gates, wrapping
    deployment sequences in \texttt{ops.change-request}, and specifying evidence requirements.
    This is a semantic upgrade; it requires author judgment and cannot be fully automated.
\end{description}

You must complete the core v1 → v2 migration before applying \texttt{gert.ops}-specific
patterns.

% ─────────────────────────────────────────────────────────────────────────────
\section{Core v1 → v2 Migration (Summary)}
\label{sec:migration-v1-to-v2}
% ─────────────────────────────────────────────────────────────────────────────

Run the automated migrator on your existing v1 runbooks:

\begin{verbatim}
$ gert migrate --to-v2 my-deploy.runbook.yaml
Migrating: my-deploy.runbook.yaml
  [REWRITE] apiVersion: runbook/v1 → runbook/v2
  [REWRITE] step type: collector → manual
  [REWRITE] step type: router → end
  [PROMOTE] meta.vars → top-level vars
  [ADD]     id field derived from meta.name
Output: my-deploy.runbook.yaml (modified in place)
Backup: my-deploy.runbook.yaml.v1.bak
\end{verbatim}

The migrator rewrites the file in place and saves a backup with the \texttt{.v1.bak} suffix.
After migration, validate the file:

\begin{verbatim}
$ gert validate my-deploy.runbook.yaml
my-deploy.runbook.yaml: valid (runbook/v2)
\end{verbatim}

For a full list of v1 → v2 breaking changes, refer to the \emph{gert v2 Design Document}, §11.

% ─────────────────────────────────────────────────────────────────────────────
\section{Adding the \texttt{gert.ops} Kit}
\label{sec:migration-add-kit}
% ─────────────────────────────────────────────────────────────────────────────

After v2 migration, add the Kit declaration and role metadata:

\begin{minted}{yaml}
# Before (v2, no kit):
apiVersion: runbook/v2
meta:
  name: deploy-service
  kind: change

# After (v2 + gert.ops):
apiVersion: runbook/v2
kit: ops/v1
meta:
  name: deploy-service
  kind: change
  roles:
    dri: "@ops-lead"
    approver: ["@change-board"]
    change-manager: "@change-board"
  sla:
    approval-timeout: 24h
    escalation-target: "@sre-lead"
\end{minted}

% ─────────────────────────────────────────────────────────────────────────────
\section{\texttt{gert.ops}-Specific Migration Patterns}
\label{sec:migration-ops-patterns}
% ─────────────────────────────────────────────────────────────────────────────

\subsection{Old \texttt{manual} → \texttt{ops.manual}}

V2 core \texttt{manual} steps that require human action and evidence should be upgraded to
\texttt{ops.manual} to gain role enforcement and structured evidence capture.

\begin{minted}{yaml}
# Before (v2 core manual):
- step:
    id: verify_health
    type: manual
    title: "Verify service health"
    instructions: |
      Check the dashboard and confirm the service is healthy.
    required_evidence:
      - kind: text
        name: health_notes

# After (ops.manual with role + typed evidence):
- step:
    id: verify_health
    type: ops.manual
    title: "Verify service health"
    requires_role: dri
    instructions: |
      Check https://grafana.internal/d/my-service.
      Confirm error rate < 0.1% and p99 latency < 200ms.
    evidence:
      - type: ops.evidence.screenshot
        label: "Health dashboard"
        required: true
        description: "Grafana dashboard showing error rate and latency."
      - type: ops.evidence.attestation
        label: "DRI health sign-off"
        required: true
        prompt: "I confirm the service is healthy."
\end{minted}

\subsection{Old Approval Steps → \texttt{ops.approval}}

V1 runbooks often used \texttt{manual} or \texttt{collector} steps as informal approval gates.
Migrate these to \texttt{ops.approval} to get timeout enforcement, formal approver lists, and
audit events.

\begin{minted}{yaml}
# Before (v1 informal approval):
- step:
    id: get_approval
    type: collector
    title: "Get manager approval"
    instructions: |
      Ask your manager to approve this change before continuing.
    required_evidence:
      - kind: text
        name: approval_note

# After (ops.approval with timeout and escalation):
- step:
    id: get_approval
    type: ops.approval
    title: "Manager approval required"
    description: |
      Deployment: {{ .service_name }} {{ .image_tag }}
      Ticket: {{ .change_ticket }}
    approvers: ["@change-board"]
    requires_any: true
    timeout: 8h
    on_timeout: escalate
    escalation_target: "@sre-lead"
    evidence:
      - type: ops.evidence.attestation
        label: "Manager approval"
        required: true
        prompt: "I approve this deployment."
\end{minted}

\subsection{Bare \texttt{cli} Deployment Steps → \texttt{ops.change-request}}

V1 deployment runbooks typically used bare \texttt{cli} steps without change management
lifecycle or rollback injection. Wrap these in \texttt{ops.change-request}.

\begin{minted}{yaml}
# Before (v1 bare cli steps):
- step:
    id: deploy
    type: cli
    command: kubectl
    args: [set, image, "deployment/api", "app={{ .image_tag }}"]

# After (wrapped in ops.change-request):
- step:
    id: execute_deployment
    type: ops.change-request
    title: "Deploy {{ .service_name }} {{ .image_tag }}"
    change-id: "{{ .change_ticket }}"
    risk: medium
    rollback:
      runbook: ./rollback.runbook.yaml
      inputs:
        service_name: "{{ .service_name }}"
        target_image_tag: "{{ .previous_image_tag }}"
    steps:
      - step:
          id: deploy
          type: ops.cli
          title: "Update deployment image"
          command: kubectl
          args: [set, image, "deployment/{{ .service_name }}", "app={{ .image_tag }}"]
          evidence:
            auto: true
\end{minted}

\subsection{Unstructured Incident Runbooks → \texttt{ops.incident}}

V1 incident runbooks were often flat sequences of manual steps with no lifecycle tracking.
Wrap the triage and mitigation steps in \texttt{ops.incident} to activate SLA timers,
commander takeover, and incident-mode evidence requirements.

\begin{minted}{yaml}
# Before (flat manual steps):
- step:
    id: triage
    type: manual
    title: "Triage the incident"
    instructions: "Check dashboards and logs."

# After (ops.incident wrapper):
- step:
    id: incident_response
    type: ops.incident
    title: "Incident response"
    severity: sev2
    incident-id: "{{ .incident_ticket }}"
    commander: "@ic-oncall"
    sla:
      mitigation-timeout: 30m
      resolution-timeout: 4h
      escalation-target: "@vp-engineering"
    steps:
      - step:
          id: triage
          type: ops.manual
          title: "Triage the incident"
          requires_role: responder
          instructions: "Check dashboards and logs."
          evidence:
            - type: ops.evidence.attestation
              label: "Triage findings"
              required: true
              prompt: "Describe the symptoms and initial hypothesis."
\end{minted}

% ─────────────────────────────────────────────────────────────────────────────
\section{Migration Validator}
\label{sec:migration-validator}
% ─────────────────────────────────────────────────────────────────────────────

After migration, use the Kit validator to check for \texttt{gert.ops}-specific issues:

\begin{verbatim}
$ gert kit validate --migration my-deploy.runbook.yaml

Validating: my-deploy.runbook.yaml (kit: ops/v1)

[WARN]  No ops.approval gate before ops.change-request wrapper (step: execute_deployment).
        Consider adding an approval gate before production changes.

[WARN]  ops.change-request at step execute_deployment has no rollback specification.
        Production changes should always declare a rollback runbook.

[ERROR] ops.manual step verify_health has no requires_role declaration.
        ops.manual steps must declare a role to enforce accountability.

[INFO]  meta.roles.dri is not set. The run initiator will become implicit DRI.
        Explicit role declaration is recommended for production runbooks.

1 error, 2 warnings, 1 info.
\end{verbatim}

The \texttt{--migration} flag enables additional warnings for patterns common in migrated
runbooks that do not yet take full advantage of \texttt{gert.ops} features.

% ─────────────────────────────────────────────────────────────────────────────
\section{Migration Validation Checklist}
\label{sec:migration-checklist}
% ─────────────────────────────────────────────────────────────────────────────

Use this checklist to verify a complete \texttt{gert.ops} migration:

\begin{center}
\small
\begin{tabular}{lp{8cm}l}
\textbf{\#} & \textbf{Check} & \textbf{Status} \\
\hline
1 & \texttt{apiVersion: runbook/v2} declared & $\square$ \\
2 & \texttt{kit: ops/v1} declared & $\square$ \\
3 & \texttt{gert validate} passes with no errors & $\square$ \\
4 & \texttt{gert kit validate} passes with no errors & $\square$ \\
5 & \texttt{meta.roles.dri} explicitly declared & $\square$ \\
6 & \texttt{meta.roles.approver} declared if approval gates exist & $\square$ \\
7 & \texttt{meta.roles.change-manager} declared if change-request exists & $\square$ \\
8 & All deployment steps wrapped in \texttt{ops.change-request} & $\square$ \\
9 & Every \texttt{ops.change-request} has a \texttt{rollback} specification & $\square$ \\
10 & Every \texttt{ops.manual} step has \texttt{requires\_role} & $\square$ \\
11 & Approval gates exist before all production-modifying sequences & $\square$ \\
12 & Evidence is captured at key decision points (not just at the end) & $\square$ \\
13 & Output redaction covers all secrets that may appear in CLI output & $\square$ \\
14 & \texttt{meta.governance.allowed\_commands} is restricted to needed tools & $\square$ \\
15 & Rollback runbook exists and has been tested in staging & $\square$ \\
\end{tabular}
\end{center}
